% Generated draft: requires mathematical and source review.
\documentclass[11pt]{article}
\usepackage[margin=25mm]{geometry}
\usepackage{fontspec}
\setmainfont{Latin Modern Roman}
\usepackage{amsmath,amssymb,mathtools,unicode-math}
\setmathfont{Latin Modern Math}
\usepackage{enumitem,needspace,xurl,hyperref}
\setlength{\parindent}{0pt}
\setlength{\parskip}{5pt}
\setlength{\emergencystretch}{3em}
\newcommand{\N}{\mathbb{N}}
\newcommand{\Z}{\mathbb{Z}}
\newcommand{\Q}{\mathbb{Q}}
\newcommand{\R}{\mathbb{R}}
\newcommand{\C}{\mathbb{C}}
\newcommand{\F}{\mathbb{F}}
\newcommand{\A}{\mathbb{A}}
\newcommand{\PP}{\mathbb{P}}
\newcommand{\E}{\mathbb{E}}
\newcommand{\eps}{\varepsilon}
\newcommand{\dd}{\,\mathrm{d}}
\begin{document}
\section*{UnsolvedMath ID: 11000233}
\textbf{Source problem code:} catalogue entry
\subsection{Definitions and mathematical statement}
Let \(W\leq\operatorname{MCG}(S)\) be the particular Whittlesey group intended by the source. The question is whether every finite-rank subgroup \(H\leq W\) is Schottky in \(\operatorname{MCG}(S)\).

In the usual mapping-class convention, a Schottky subgroup is a finitely generated free convex cocompact subgroup: its orbit in Teichm\"uller space is quasiconvex in the intended metric. If \(W\) is being used as a free group, finite rank means finitely many free generators. The construction of \(W\), surface \(S\), exact convex-cocompactness definition and treatment of rank-zero/rank-one groups are not supplied; they must be fixed to interpret the universal quantifier. No passage to powers of generators or to a smaller subgroup is allowed as a substitute for \(H\) itself being Schottky, and no uniform quasiconvexity constant over all \(H\) is requested.

The continuation gives a consequence of its Theorem 2.1: if \(H\) has a finite-index Schottky subgroup, then the associated extension \(\Gamma_H\) has a finite-index word-hyperbolic subgroup, hence is itself word hyperbolic. The usual extension has the form
\[
 1\longrightarrow\pi_1(S)\longrightarrow\Gamma_H\longrightarrow H\longrightarrow1.
\]
Its precise construction and Theorem 2.1 hypotheses, especially the surface convention, are missing. This is reported background about virtually Schottky groups, not the stronger answer that all finite-rank subgroups of \(W\) are Schottky.
\subsection{Short English statement}
Is every finite-rank subgroup of the specified Whittlesey group already a free convex cocompact mapping-class subgroup, with the associated virtual-Schottky extension consequence kept separate?
\subsection{Sources}
\begin{itemize}
\item \emph{Problems on Mapping Class Groups and Related Topics}, edited by Farb, Problem 3.5 and continuation, supplied PDF page 272. \url{https://www.math.uchicago.edu/~farb/papers/mcgbook.pdf}. Recorded access 2026-07-29; extraction pdf-font-structure.
\item Whittlesey's group and the source Theorem 2.1 are named without their definitions, full statements or separate bibliography entries in this record.
\item Local UnsolvedMath registry, \texttt{data/unsolvedmath/problems.json}, permanent ID \(11000233\), exact source code \texttt{catalogue entry}; supplied OPEN-TRIAGE summary reports unverified status.
\end{itemize}
\end{document}
